\section{Related Work}
\label{sec:related}

\para{Internal-state deception detection.}
The dominant approach to detecting LLM deception operates on model internals. \citet{azaria2023internal} train \saplma, a classifier on hidden-layer activations, achieving 71--83\% accuracy at distinguishing true from false statements across multiple LLMs. \citet{burns2023discovering} discover latent knowledge in LLMs without supervision by identifying directions in activation space that satisfy logical consistency properties. \citet{marks2025probing} extend this line by probing the geometry of truth representations, finding structured truth directions that generalize across logical transformations. \citet{levinstein2024truth} demonstrate that these truth directions are universal, transferring across topics, datasets, and even models.

More recently, \citet{long2025truthful} use sparse autoencoders to show that deceptive instructions induce specific representational shifts: truthful and deceptive representations occupy distinct subspaces, with identifiable features that ``flip'' under deceptive prompts. \citet{huan2025can} apply mechanistic interpretability to locate sparse lying circuits concentrated in specific attention heads, and extract steering vectors that can control deceptive behavior. These findings collectively establish that LLMs encode truthfulness in their internal representations. Our work asks the complementary question: do these internal differences manifest in the output text?

\para{Deception in LLMs.}
\citet{jones2024lies} provide a taxonomy of LLM deception, distinguishing influence, persuasion, manipulation, and strategic deception. They identify three mechanisms by which LLMs produce deceptive behavior: pre-training data patterns, fine-tuning rewards, and in-context instruction following. \citet{hubinger2024sleeper} demonstrate that deceptive behavior can persist through safety training, with ``sleeper agents'' maintaining hidden objectives despite RLHF and adversarial training. \citet{burger2025benchmarking} introduce metrics for rigorously evaluating deception probes, testing whether they detect genuine deception or rely on superficial correlations.

\para{Surface-level text analysis.}
Comparatively little work examines whether deceptive LLM output is distinguishable from truthful output at the text level. \citet{schuster2019limitations} find that stylometric features have limited ability to distinguish machine-generated fake news from real news, especially under style transfer attacks. \citet{chen2023fake} observe that fake news detectors classify LLM-generated text as fake regardless of veracity, suggesting detectable stylistic properties in LLM text. However, neither study examines paired truthful vs.\ deceptive outputs from the same model under controlled conditions---the setting we investigate here.

\para{Positioning.}
Our work fills a gap between internal-representation approaches and practical black-box detection. While prior work establishes that truthful and deceptive processing differ \emph{inside} the model, we test whether these differences propagate to the \emph{output text}. We use controlled paired generation with multiple deception strategies, extract simple linguistic features, and evaluate detection accuracy---providing a practical complement to the mechanistic interpretability literature.
